\section{Contexto y Motivación}
\label{sec:contexto}

La gestión financiera personal representa uno de los desafíos más significativos de hoy en día. El volumen de transacciones digitales crece exponencialmente \cite{fintech2023}, mientras que los usuarios se enfrentan diariamente a la necesidad de comprender sus patrones de gasto y establecer objetivos de ahorro realistas. Sin embargo, las herramientas disponibles ofrecen o bien análisis retrospectivo básico o requieren conocimientos técnicos avanzados para su configuración y uso efectivo.

\subsection{Problemática Actual}
\label{subsec:problematica}

El panorama actual de las aplicaciones de gestión financiera personal presenta limitaciones significativas en tres dimensiones fundamentales.

\textbf{Limitaciones de interpretación:} Las soluciones comerciales ofrecen dashboards genéricos que presentan datos de forma visualmente atractiva pero carecen de capacidad de interpretación contextual \cite{conversational_ai2023}. Estas herramientas pueden mostrar que un usuario ha gastado un 20\% más en restaurantes, pero son incapaces de explicar por qué ha ocurrido este aumento, qué tendencia histórica lo respalda o cómo podría optimizarse.

\textbf{Barreras de usabilidad:} Las soluciones basadas en hojas de cálculo y software de contabilidad requieren un alto grado de conocimiento técnico. Esta barrera limita la adopción de prácticas financieras saludables y perpetúa la toma de decisiones basada en intuición en lugar de análisis de datos objetivos \cite{personal_finance_ai2024}.

\textbf{Ausencia de continuidad:} Los sistemas existentes carecen de memoria entre sesiones. Cada interacción parte de cero, lo que impide que el sistema aprenda las preferencias del usuario, recuerde objetivos pactados o detecte cambios significativos en los patrones de gasto a lo largo del tiempo.

\subsection{Fundamentos Tecnológicos}
\label{subsec:fundamentos}

La reciente evolución de los modelos de lenguaje grandes (LLMs) y las arquitecturas de agentes basados en grafos ha abierto nuevas posibilidades en el campo de la asistencia financiera personal \cite{tool_using_agents2023}. Estos sistemas combinan la capacidad de comprensión del lenguaje natural con el razonamiento estructurado, permitiendo crear interfaces conversacionales que pueden entender consultas complejas y realizar análisis multi-paso.

La arquitectura LangGraph \cite{langgraph2024} ofrece un marco ideal para este tipo de aplicaciones al permitir el diseño de agentes que deciden dinámicamente qué herramientas utilizar, encadenan múltiples operaciones y mantienen estado entre interacciones. A diferencia de los sistemas basados en pipelines fijos, los agentes LangGraph adaptan su comportamiento a las necesidades específicas de cada consulta mediante un grafo de nodos y aristas que define el flujo de control.

\subsection{Motivación del Proyecto}
\label{subsec:motivacion}

Este proyecto surge de la confluencia de dos motivaciones complementarias. Por un lado, la necesidad técnica de explorar las capacidades de los agentes con herramientas en un dominio real y bien acotado. Por otro, el objetivo práctico de crear una herramienta útil para el usuario final que supere las limitaciones de los enfoques actuales.
